\section{Supplementary material}\label{supplementary-material}

\begin{itemize}
\item \textbf{S1.} Circuit schematic of the FLEX-DBS board (PDF).
\item \textbf{S2.} Firmware for the BGM220S (Silicon Labs Simplicity Studio project), v0.5. Repository: \url{https://github.com/Neurotech-Hub/Creed-DBS-SiLabs}
\item \textbf{S3.} MouseCap iOS application source. Repository: \url{https://github.com/Neurotech-Hub/MouseCap-iOS}
\item \textbf{S4.} BLE protocol specification (\texttt{MCU\_BLE\_PROTOCOL.md} in S3): service and characteristic UUIDs, command grammar, parameter keys and ranges, and response format.
\item \textbf{S5.} Shroud design (Fusion 360): \url{https://a360.co/3TqHAxf}
\item \textbf{S6.} Electrode holder design (Fusion 360): \url{https://a360.co/4rEBmq2}
\item \textbf{S7.} Extended comparison table (Table S7, below): programmable ranges, dimensions and release status for the systems in Table 4.
\end{itemize}

\begin{landscape}
\textbf{Table S7.} Extended comparison of FLEX-DBS with published wireless rodent stimulators (full version of Table 4). Values for published systems are as reported by the authors; blank cells, and sizes omitted from the mass column, are not reported in the source. Where a source gives two values, the second is shown in parentheses with its location. CC, constant current; CV, constant voltage.

{\scriptsize
\setlength{\tabcolsep}{2pt}
\begin{longtable}{@{}>{\raggedright\arraybackslash{}}p{0.09\linewidth}>{\raggedright\arraybackslash{}}p{0.07\linewidth}>{\raggedright\arraybackslash{}}p{0.08\linewidth}>{\raggedright\arraybackslash{}}p{0.08\linewidth}>{\raggedright\arraybackslash{}}p{0.05\linewidth}>{\raggedright\arraybackslash{}}p{0.10\linewidth}>{\raggedright\arraybackslash{}}p{0.14\linewidth}>{\raggedright\arraybackslash{}}p{0.12\linewidth}>{\raggedright\arraybackslash{}}p{0.12\linewidth}>{\raggedright\arraybackslash{}}p{0.09\linewidth}@{}}
\toprule
System & Species tested & Mass with battery; size & Channels & Output & Supply / compliance & Amplitude, pulse width, frequency; waveform & Parameter control & Current draw; operating duration & Open release \\
\midrule
\endhead
\textbf{FLEX-DBS (this work)} & Mouse & 3.71 g with shroud (two zinc-air cells); 19.5 × 22.5 × 8.72 mm & \textbf{2; shared amplitude and timing}$^{\dagger}$ & CC & ±5 V rails; $\sim$4.6 V swing (measured) & 0--625 µA ideal (compliance-limited in vivo); 90--600 µs set; $\sim$83--167 Hz delivered; asymmetric biphasic, 1:5 recovery; \textbf{continuous or burst}$^{\dagger}$ & \textbf{Real-time BLE from an iOS app; autonomous after disconnect}$^{\dagger}$ & 3.15--3.42 mA stimulating; $\sim$3 d per cell pair (projected) & Schematic, firmware, app, protocol (S1--S4) \\
t-IPG, \citet{paulat2026} & Mouse (s-IPG also in rats and mice) & 0.9 g (s-IPG 1.6 g); 5.3 mm thick & 2, time-delayed & CC & 3 V lithium cell, direct & 1--100 µA, 50--700 µs, 2--257 Hz; cathodic phase with long low-amplitude anodic phase; theta- and beta-burst modes & Predefined paradigms via near-field transponder & 21.3 µA maximum; 49 d voltage-compliant & Partial (GitHub) \\
STELLA, \citet{plocksties2021} & Rat (proposed for mice $>$20 g) & 4.0 g encapsulated (CR1225), Ø17.0 × 9.0 mm; 2.6 g (CR1216), Ø15.5 × 7.5 mm & 2, synchronized & CC & Boost 3.1--3.7 V; 2.2--2.8 V at electrode & 42--100 µA, 1 µs--100 \% duty, 10--5000 Hz; monophasic, passive charge balance & Firmware set before implantation; Hall sensor steps through presets & 7.6--8.3 µA; $\ge$12 wk in vitro, 42 d in vivo & Yes, open license (GitHub) \\
\citet{grotemeyer2024} & Mouse & 2 g; 18 × 16 × 5 mm with housing & 1 & CC & 3 V (two 1.5 V cells) & 10--500 µA, 60--500 µs, 10--200 Hz; monophasic, passive discharge & Cable programmer, reprogrammable after implantation; reed switch on/off & $\sim$20 µA; 52 d in vitro & Analysis code only; commercialized \\
\citet{fleischer2020} & Mouse & 2.8 g; board 9 × 9 × 1 mm (abstract: 10 × 10 mm) & 1 & CC & 3 V (two V364 cells) & 0--500 µA (text: up to 300 µA), 60--500 µs, 10--200 Hz (text: 10--500 Hz); monophasic, passive discharge & Programmed before mounting; switch on/off & $\sim$20 µA; $\sim$30 d in vitro & \\
\citet{dehaas2012} & Mouse & 2.1 g with batteries and lead; Ø8 × 30 mm & 2 (bilateral), pulses alternate between channels & CC (adaptive series resistance) & 4.65 V (three 337 silver-oxide cells); regulates into 15--45 kΩ & 20--100 µA, 60 µs, 131 Hz fixed; symmetric biphasic & Amplitude set by reprogramming in a fixture (code table and resistor changes); magnet on/off; IR LED status & 15 d standby, 10 h stimulating & \\
\citet{alpaugh2019} & Mouse & 4.7 g & 1 & CC & 3 V (two 393 cells); 2.1 V output headroom & Fixed 149 µA, $\sim$103 µs, $\sim$130 Hz; balanced biphasic & Preprogrammed via external circuit; no radio & $\sim$2 wk per cell set, $\ge$30 d with weekly exchange & Component list only \\
\citet{kouzani2013} & Rat (designed for mice) & 3.27 g (head-mount) & 1 & CC & 3.1 V (two silver-oxide cells) & Fixed 200 µA, 90 µs, 130 Hz; monophasic, capacitor-coupled & Preprogrammed; changes require reflashing or a resistor change & $\sim$870 µA; 12.3 d on bench, 7 d in vivo & Schematic and pseudocode \\
\citet{pinnell2018} & Rat & 2.8 g with battery and head cap (text: 2.7 g); device Ø12.5 × 5 mm & 2 & CC & 12.29 V regulated & 20 µA--2 mA, 10 µs--100 \% duty, 0.1--5000 Hz; monophasic or biphasic; dual-frequency burst & Wired programmer; amplitude by potentiometer; magnet on/off & 35 µA in low-power state; $\sim$30 h stimulating & Design files and assembly instructions (Figshare) \\
\citet{burton2021} & Rat & 0.046 g, battery-free & 1 bipolar electrode & CV & 1.5--5.5 V (regulated) & 1.5--5.5 V, $\ge$20 µs pulses, up to 25 kHz; biphasic & Real time via modulation of RF power & 9.35 mW while stimulating; $>$36 d in vivo with RF power & \\
\citet{shen2026} & Rat & 12.5 g (7.5 g battery); 34 × 42 × 15 mm with housing & Not stated; used unilaterally & CV & 3.7 V Li-ion & 2.3--3 V, 40--340 µs, 80--150 Hz; biphasic & Real-time BLE (MATLAB interface) & $\ge$300 h on 400 mAh & \\
\bottomrule
\end{longtable}
$^{\dagger}$New paradigm enabling.
}
\end{landscape}
